% QCM de sortie — IMSV séance 4 (Vecteurs et systèmes de coordonnées).
% Compiler avec LuaLaTeX (en-tête XML en UTF-8) :  lualatex qcm-sortie-seance04.tex
% Produit qcm-sortie-seance04-moodle.xml, à importer dans Moodle :
%   Banque de questions > Importer > format « XML Moodle ».
% Notation du cours (M. Randour, IMSV-4). Maths rendues par MathJax dans Moodle.
\documentclass[11pt]{article}
\usepackage{fontspec}
\usepackage{amsmath,amssymb}
\usepackage[a4paper,margin=2cm]{geometry}
\usepackage{moodle}

\begin{document}

\begin{quiz}{IMSV/Séance 4 — QCM de sortie}

% ---------------------------------------------------------------- Q1
\begin{multi}[penalty=0, shuffle=true, numbering=abc,
  feedback={Le vecteur $\overrightarrow{AB}$ décrit le déplacement de $A$ vers $B$ :
  $\overrightarrow{AB} = (x_B - x_A, y_B - y_A)$. Sa norme se calcule par Pythagore.
  Voir \emph{Coordonnées cartésiennes} et \emph{Norme d'un vecteur}.}]{S4-Q1 Coordonnées et norme d'un vecteur}
Soient les points $A = (1, -2)$ et $B = (4, 2)$. Quelles sont les coordonnées du vecteur
$\overrightarrow{AB}$ et sa norme ?
\item[fraction=100, feedback={Exact : extrémité moins origine, $(4 - 1, 2 - (-2)) = (3, 4)$, puis
  Pythagore : $\sqrt{9 + 16} = 5$.}]
  $\overrightarrow{AB} = (3, 4)$ et $\Vert\overrightarrow{AB}\Vert = 5$
\item[feedback={Tu as calculé origine moins extrémité : c'est $\overrightarrow{BA} = -\overrightarrow{AB}$,
  le vecteur de sens opposé.}]
  $\overrightarrow{AB} = (-3, -4)$ et $\Vert\overrightarrow{AB}\Vert = 5$
\item[feedback={Tu as additionné les coordonnées des points. Les coordonnées d'un vecteur s'obtiennent
  par « extrémité moins origine ».}]
  $\overrightarrow{AB} = (5, 0)$ et $\Vert\overrightarrow{AB}\Vert = 5$
\item[feedback={Les coordonnées sont justes, mais la norme n'est pas la somme des coordonnées :
  c'est $\sqrt{3^2 + 4^2}$ (Pythagore).}]
  $\overrightarrow{AB} = (3, 4)$ et $\Vert\overrightarrow{AB}\Vert = 7$
\end{multi}

% ---------------------------------------------------------------- Q2
\begin{multi}[penalty=0, shuffle=true, numbering=abc,
  feedback={Pour $\vec v$ de norme $r$ formant l'angle $\theta$ avec l'horizontale :
  $\vec v = (r\cos\theta, r\sin\theta)$. Les signes des coordonnées doivent correspondre au
  quadrant. Voir \emph{De la norme aux coordonnées}.}]{S4-Q2 De la norme aux coordonnées}
Un vecteur $\vec v$ a pour norme $4$ et fait un angle de $120^\circ$ avec l'horizontale (sens
trigonométrique). Quelles sont ses coordonnées ?
\item[fraction=100, feedback={Exact : $\vec v = 4(\cos 120^\circ, \sin 120^\circ)
  = 4\left(-\frac12, \frac{\sqrt3}{2}\right)$. Deuxième quadrant : $v_x < 0$ et $v_y > 0$.}]
  $\left(-2, 2\sqrt3\right)$
\item[feedback={Ce sont les coordonnées pour $60^\circ$. À $120^\circ$, on est dans le deuxième
  quadrant : le cosinus est négatif.}]
  $\left(2, 2\sqrt3\right)$
\item[feedback={Tu as échangé cosinus et sinus (c'est l'angle de $150^\circ$). La première coordonnée
  est $r\cos\theta$, la seconde $r\sin\theta$.}]
  $\left(-2\sqrt3, 2\right)$
\item[feedback={C'est le point du cercle trigonométrique (norme $1$). Il faut encore multiplier par la
  norme $4$.}]
  $\left(-\frac12, \frac{\sqrt3}{2}\right)$
\end{multi}

% ---------------------------------------------------------------- Q3
\begin{multi}[penalty=0, shuffle=true, numbering=abc,
  feedback={$\vec u$ et $\vec v$ (non nuls) sont colinéaires s'il existe un réel $\lambda$ tel que
  $\vec v = \lambda \cdot \vec u$, le \textbf{même} $\lambda$ pour toutes les coordonnées.
  Voir \emph{Colinéarité}.}]{S4-Q3 Colinéarité}
Lequel de ces vecteurs est colinéaire à $\vec u = (2, -3)$ ?
\item[fraction=100, feedback={Exact : $(-4, 6) = -2 \cdot (2, -3)$. Même direction, sens opposé.}]
  $(-4, 6)$
\item[feedback={Échanger les coordonnées ne donne pas un multiple : il faudrait $\lambda = \frac32$
  pour la première coordonnée et $\lambda = \frac23$ pour la seconde.}]
  $(3, -2)$
\item[feedback={La première coordonnée demande $\lambda = 2$, la seconde $\lambda = -2$ : pas de
  $\lambda$ commun, donc pas colinéaires.}]
  $(4, 6)$
\item[feedback={Il faudrait $\lambda = -1$ pour la première coordonnée et $\lambda = 1$ pour la
  seconde : pas colinéaires.}]
  $(-2, -3)$
\end{multi}

% ---------------------------------------------------------------- Q4
\begin{multi}[penalty=0, shuffle=true, numbering=abc,
  feedback={On calcule le produit scalaire en coordonnées, puis on l'égale à
  $\Vert\vec u\Vert \Vert\vec v\Vert \cos\theta$ avec $\theta \in [0, \pi]$. Le signe de
  $\vec u \cdot \vec v$ annonce déjà si l'angle est aigu, droit ou obtus.
  Voir \emph{Angles via le produit scalaire}.}]{S4-Q4 Angle via le produit scalaire}
Quel est l'angle entre $\vec u = (2, 0)$ et $\vec v = (1, 1)$ ?
\item[fraction=100, feedback={Exact : $\vec u \cdot \vec v = 2$, $\Vert\vec u\Vert = 2$,
  $\Vert\vec v\Vert = \sqrt2$, donc $\cos\theta = \frac{2}{2\sqrt2} = \frac{\sqrt2}{2}$.}]
  $\frac{\pi}{4}$ (soit $45^\circ$)
\item[feedback={$\cos\frac{\pi}{3} = \frac12$. Recalcule
  $\cos\theta = \frac{\vec u \cdot \vec v}{\Vert\vec u\Vert \Vert\vec v\Vert}$ : on obtient
  $\frac{\sqrt2}{2}$.}]
  $\frac{\pi}{3}$ (soit $60^\circ$)
\item[feedback={Le produit scalaire vaut $2 > 0$ : l'angle est aigu, il ne peut pas être obtus.}]
  $\frac{3\pi}{4}$ (soit $135^\circ$)
\item[feedback={Des vecteurs orthogonaux auraient un produit scalaire nul ; ici
  $\vec u \cdot \vec v = 2$.}]
  $\frac{\pi}{2}$ (soit $90^\circ$)
\end{multi}

% ---------------------------------------------------------------- Q5
\begin{multi}[penalty=0, shuffle=true, numbering=abc,
  feedback={$\vec u \times \vec v$ est perpendiculaire au plan de $\vec u$ et $\vec v$, sa norme est
  l'aire du parallélogramme construit sur eux ($\Vert\vec u\Vert \Vert\vec v\Vert \sin\theta$) et son
  sens est donné par la règle de la main droite ; $\vec v \times \vec u = -(\vec u \times \vec v)$.
  Voir \emph{Produit vectoriel dans $\mathbb{R}^3$}.}]{S4-Q5 Produit vectoriel}
Dans $\mathbb{R}^3$, soient $\vec u = (2, 0, 0)$ et $\vec v = (0, 3, 0)$. Que vaut
$\vec u \times \vec v$ ?
\item[fraction=100, feedback={Exact : perpendiculaire au plan $(x, y)$, donc porté par l'axe des
  $z$ ; norme égale à l'aire du rectangle, $2 \cdot 3 = 6$ ; main droite (index sur $\vec u$, majeur
  sur $\vec v$) : vers les $z$ positifs.}]
  $(0, 0, 6)$
\item[feedback={La direction et la norme sont justes, mais le sens est inversé : c'est
  $\vec v \times \vec u$. Applique la règle de la main droite dans l'ordre $\vec u$, puis $\vec v$.}]
  $(0, 0, -6)$
\item[feedback={C'est $\vec u + \vec v$. Le produit vectoriel est perpendiculaire au plan de $\vec u$
  et $\vec v$.}]
  $(2, 3, 0)$
\item[feedback={C'est le produit scalaire (les vecteurs sont orthogonaux). Le produit vectoriel est un
  \textbf{vecteur}, de norme égale à l'aire du parallélogramme construit sur $\vec u$ et $\vec v$.}]
  $0$
\end{multi}

\end{quiz}
\end{document}
